\section{Positive pressures with shared boundaries}\label{ms:pressure}
The estimates just proved have different complete tensors on their
right sides. This section compares their entire contractions. The
comparison does not assert an isotopy of the corresponding tangles.

Let $\mathcal G$ be a finite tree of bare junctions, each of full
valence at least three. Its internal edges are physical turns and
its ordinary external profiles are $g+c_j$. A covariant junction
label is the index inserted in its trace; its rooted ordinary output
is the complementary label. Thus a physical edge identifies equal
covariant labels, while a hidden unturned resolution edge identifies
complementary ones. The cubic action is symmetric:
\begin{equation}\label{ms:cubic-action}
 t(x,y,z)=m(x,y,1-z)
 =J(x+y)+J(x+z)+J(y+z)-J(x)-J(y)-J(z)-J(x+y+z-1).
\end{equation}
Choose an ordered binary resolution at each junction. Let
$V_D(a_1,\ldots,a_D)$ denote the maximum of its gate action with
the covariant boundary labels fixed, and let $F(\mathcal G)$ be
the maximum after including all ordinary profiles and shared edges.

\begin{lemma}[Strict resolutions and constrained maxima]
\label{ms:pressure-existence}
The strict support of a whole junction is
\begin{equation}\label{ms:whole-support}
 a_i>0,\qquad a_i+a_j<1\ (i\ne j),\qquad \sum_i a_i>1.
\end{equation}
Every prescribed binary resolution realizes that support. The tree
pressure has a unique strict maximum. The same assertion holds if
at most one external incidence of each whole junction is fixed at
$1/2$, including small independent variations of those fixed labels.
The constrained pressure is analytic there with negative definite
Hessian in its free boundary labels.
\end{lemma}
\begin{proof}
For a product of inputs $a_i$, its strict output interval is
\[
 \max_i a_i<b<\min(1,\sum_i a_i),\qquad a_i+a_j<1.
\]
One proves this for any binary association by induction. A composite
child has output interval $(L_S,U_S)$ with
$L_S=\max_{i\in S}a_i$, $U_S=\min(1,\sum_{i\in S}a_i)$;
a singleton is fixed. Choose the two child outputs below $b$ with
sum between $b$ and one. Their lower sum is below one by the pair
condition and their upper sum exceeds $b$. The singleton cases use
the same inequalities directly. Conversely the gate at which two
inputs first meet forces their pair sum below one. Closing the last
incidence makes $b=1-a_D$ and gives \eqref{ms:whole-support}.
Its closed version follows by the same finite interval rules.

Assign $3/8$ to every physical edge and ordinary boundary. This is
a strict witness at every junction. If one boundary at a junction
is fixed at half, its largest pair is $7/8$ and its total is at least
$5/4$; otherwise these bounds are $3/4$ and $9/8$. The witness
persists under small independent variations of the fixed halves.

Here the endpoint argument must also retain those fixed boundaries.
For positive triples of sum one put
$\mathcal V(a,b,c)=-\pi(J(a)+J(b)+J(c))$. Its Hessian is
negative definite on the angle plane: the diagonal cotangent sums
are positive and
$\cot\pi a\cot\pi b+\cot\pi b\cot\pi c+
\cot\pi c\cot\pi a=1$.
If $p=z-x$, $q=z-y$, $r=x+y-z$, $w=1-x-y$, then
\begin{equation}\label{ms:gate-atoms}
 \pi m(x,y,z)=\mathcal V(y,x,w)
   +\mathcal V(p,q,2r+w)+\mathcal V(p+q,r,r+w).
\end{equation}
The first atom's zero quadratic form forces $dx=dy=0$ and the
second then forces $dz=0$. Thus a strict gate has a negative
definite Hessian in all three variables. Also
$g'(a)=2\log\cot\pi a$ and $g''<0$ on $(0,1/2)$.

At a nonvertex flat triple $(0,a,1-a)$ an inward segment has
derivative $+\infty$, because $J(tu)=tu\log t+O(t)$.
At a vertex flat triple the $t\log t$ terms cancel, so no negative
infinite term can cancel this barrier. Root the resolved tree at an
output or a fixed half incidence. Output zero forces both product
inputs to zero. Output one either has two interior inputs summing to
one, which supplies a nonvertex flat atom, or propagates an endpoint
to a child. A physical turn exchanges zero and one; an unturned
resolution edge preserves the ordinary endpoint. Following this
descent in a finite tree reaches an ordinary leaf endpoint, where
the inward derivative of $g$ is infinite, or an interior pair summing
to one. It cannot terminate at an interior fixed boundary. No
maximizer therefore has an endpoint wire. With all wires interior,
any remaining equality in the gate support supplies a nonvertex
flat atom in \eqref{ms:gate-atoms}. The maximum is strict.

Every free coordinate appears in a strict gate or a leaf term, so
the summed Hessian is negative definite. Concavity gives uniqueness.
Including the variable prescribed boundaries in that quadratic form
and taking a Schur complement proves the asserted strict concavity
of their pressure. The implicit function theorem and compactness
give the analytic local maximizing branch.
\end{proof}

\begin{lemma}[Reflection and deletion of turns]\label{ms:deletion}
An unturned rooted side with ordinary boundaries has profile $f$ on
$(0,1)$ satisfying $f(b)\ge f(1-b)$ for $b\ge1/2$.
The inequality is strict above half unless the entire rooted side
is one trivalent junction with two ordinary leaves. Consequently
deleting any set $S$ of physical turns and fusing their adjacent
ordered junctions gives
\begin{equation}\label{ms:pressure-deletion}
 F(\mathcal G/S)\le F(\mathcal G).
\end{equation}
Every internal physical label at the original maximum is at most
half. Equality occurs only for the one internal edge of the entire
two-trivalent-junction tree with two ordinary leaves on each side.
\end{lemma}
\begin{proof}
We use the stationary response equations of
\cite[Lemma 4.1 and Proposition 4.2]{Profiles}, recording their
inequality consequence to preserve the constants. Put
\[
 B=\cot\pi b,\quad \kappa=e^{-f'(b)},\quad
 X^2=\frac{1+B^2}{1+\kappa},\quad Y^2=\kappa X^2.
\]
An ordinary leaf has $(B,X,Y)=(B,B,1)$. A composite has
$-Y<B<X$. At a product the inputs and output have common $X$ and
\begin{equation}\label{ms:response}
 B_o=\frac{B_1B_2+X^2-1}{B_1+B_2},\qquad
 Y_o=\frac{Y_1Y_2}{B_1+B_2},\qquad B_1+B_2>0.
\end{equation}
A physical turn sends $(B,X,Y)$ to $(-B,Y,X)$. Proper responses
remain proper. To check the possible signs directly, for $X\ge1$
the input inequalities force $B_i>0$ and
$B_o-X=((X-B_1)(X-B_2)-1)/(B_1+B_2)<0$.
For $X<1$, write $B_i=\sqrt{1-X^2}\cot\alpha_i$,
$Y_i=\sqrt{1-X^2}\csc\alpha_i$; the denominator gives
$\alpha_1+\alpha_2<\pi$ and the product response has this angle.
The asserted inequalities follow. Thus these formulas also apply
to every turned composite side in the present tree.

At the last gate of an unturned product,
\[
 Y^2=(B_1-B)(B_2-B)\le(X-B)^2,
 \qquad X\ge X_2(B):=\frac{B+\sqrt{B^2+2}}2.
\]
The factors are positive by strict support. Since
$\log(X^2/(1+B^2-X^2))$ increases with $X$,
\begin{equation}\label{ms:profile-derivative}
 f'(b)\ge f_2'(b)=2\log(2X_2(B)^2),
\end{equation}
strictly when either input is composite. The actual two-leaf
profile $f_2$ is symmetric, because $X_2(B)X_2(-B)=1/2$.
Integrating \eqref{ms:profile-derivative} from $1-b$ to $b$
proves the reflection statement, with all additive constants
unchanged. The exceptional side is precisely the two ordinary
leaves. At $b=1/2$ its derivative is zero; every other side has
positive derivative there.

Across an original edge the two ordinary outputs are both $1-z$,
so its cut pressure is $f_A(1-z)+f_B(1-z)$. Above half it
decreases strictly; at half it decreases unless both sides are
exceptional. This proves the internal-label statement. The
boundary labels themselves remain strictly between zero and half.

Exact fusion is supplied by the full metric:
\begin{equation}\label{ms:fusion}
 \sum_{i=0}^d\phi_0(Ae_i)\phi_0(e_{d-i}B)=\phi_0(AB).
\end{equation}
Indeed expand $A=\sum_i\phi_0(Ae_i)e_{d-i}$. Reflecting one
endpoint of an equal-index physical edge gives exactly this
complementary pairing. The induced order is the boundary order of
the two vertex disks and their joining band; it substitutes the
child product in its original position. No attached branches are
permuted. The pressures before and after this fusion are
\[
 F_T=\max_b(f_A(b)+f_B(b)),\qquad
 F_I=\max_b(f_A(b)+f_B(1-b)).
\]
For $b\ge1/2$ use reflection on $B$, and for $b\le1/2$ use it
on $A$ at $1-b$. Thus $F_I\le F_T$. A fused junction has valence
$D_1+D_2-2\ge4$, retains ordinary boundary constants, and remains
in the same class. Apply this argument successively to the actual
profiles of each current fused tree to obtain
\eqref{ms:pressure-deletion}. Equality is possible; no strict
source loss has yet been inferred.
\end{proof}

\begin{lemma}[Several high boundaries]\label{ms:half-boundaries}
Fix marked covariant boundary labels $z_i\ge1/2$, at most one at
each retained whole junction, without leaf actions on those marks.
Let $\Psi(\mathbf z)$ be the maximum of the same complete action.
Then $\Psi(\mathbf z)\le\Psi(\tfrac12,\ldots,\tfrac12)$.
Moreover, if a whole junction has $k\ge2$ half incidences, then
\begin{equation}\label{ms:cardinal-action}
 V_D(\underbrace{\tfrac12,\ldots,\tfrac12}_{k},x_1,\ldots,x_{D-k})
       =\sum_j g(x_j),\qquad 0\le x_j\le\tfrac12.
\end{equation}
\end{lemma}
\begin{proof}
Lemma~\ref{ms:pressure-existence} gives a strict analytic constrained
maximum at the all-half point. Root there at mark $i$. Each other
fixed half terminal has $B=0$ and the positive response
$\kappa=e^{m_x}$ supplied by its incident strict gate; hence it
is proper, $-Y<B<X$. This response is not a new leaf action.
Propagate it by \eqref{ms:response}. The final product has ordinary
output half and two inputs below half, so $B_1,B_2>0$ and
\[
 Y^2=1-X^2=B_1B_2\le X^2,
 \qquad \partial_{z_i}\Psi=-\log(X^2/Y^2)\le0.
\]
The symmetric cubic action \eqref{ms:cubic-action} shows that
rerooting changes no constrained pressure. Its joint concavity
therefore gives the claimed simultaneous bound by the supporting
plane at the all-half point. Closed boundary states follow by
continuity. No assertion of coordinate monotonicity with other
marks already above half is needed.

For the second assertion use odd levels $N=2m+1$ solely to
calibrate the real action. The polynomial $e_m^2$ is the full
Hermite projector at the simple node $-2$: it vanishes with both
jets at every double node, and
\[
 e_m(-2)=1,\qquad\phi_0(e_m^2)=1,\qquad
 \phi_0(e_m^2F)=F(-2).
\]
For example $S_m^2=N$ and $P_0'(-2)=N^2$ give the first value;
the top coefficient of $e_m^2$ gives the second. Hence
$\phi_0(e_m^k\prod e_{n_j})=\prod e_{n_j}(-2)$ for $k\ge2$.
Take its uniform closed factorial rate, using the positive ordered
resolution \eqref{ms:gate}. Here is a support-preserving approximation,
including the saturated pair faces. Keep every half index exactly
$m=d/2$ and set the other inputs to $n_j=\lfloor d x_j\rfloor$.
All pair sums are at most $d$, and the two retained halves make
the total at least $d$. The integer interval rules of the complete
product therefore give at least one supported resolution.
In a sequential resolution its running outputs obey only integer
difference inequalities: $r_a-r_{a-1}\in[0,n_a]$,
$r_a\ge n_a$, $r_{a-1}\le d-n_a$, with the fixed first and last
outputs. Encode bounds by a further coordinate $r_0=0$.

For an arbitrary real supported resolution let $Y_a$ be its
outputs multiplied by $d$, and $Y_0=0$. Every rounded difference
inequality $r_j-r_i\le b_{ij}$ is violated by $Y$ by at most two.
The integer system is feasible, so its directed constraint graph
has no negative cycle. Let $d_{ij}$ be its shortest-path lengths;
all coordinates are bounded, so these lengths are finite. A
shortest path can be taken simple, with at most $D$ edges, giving
$d_{ij}\ge Y_j-Y_i-2D$. The integer vector
$z_j=\min_i(\lfloor Y_i\rfloor+d_{ij})$ satisfies every difference
inequality and $Y_j-2D-1\le z_j\le Y_j$. Subtracting $z_0$ fixes
the zero coordinate without changing any inequality. The resulting
supported integer outputs differ from $Y$ by at most $2D+1$ and
satisfy both prescribed endpoint outputs exactly. Thus every real
resolved state has supported lattice approximants within $O(1/d)$.
Conversely every lattice state is a point of the same limiting
closed support. Uniform closed factorial estimates give $O(\log N)$
errors for these bounded displacements. Positive state counting
then identifies the exact finite rate with the maximum action.
This proves
\eqref{ms:cardinal-action} in its relative interior, and continuity
adds zero arms and further half arms. The real action is independent
of parity; no even-level simple-node projector is asserted.
\end{proof}

These two lemmas give a useful reassembly principle. Remove a
nonempty collection of pairwise nonadjacent moving junctions, each
with at least two high incidences, and bound each by the sum of
$g$ on its low incidences and zero on the high ones. The neighboring
retained complete hubs have at most one high incidence by their
pairwise support. Mark those high boundaries and charge a single
$g$ to each low arm of a removed hub. Optimize each coupled old
region, move all its marks to half by Lemma~\ref{ms:half-boundaries},
then reinsert the removed hubs by \eqref{ms:cardinal-action}.
Their cardinal action consumes the charged $g$ profiles exactly
once. Original low external profiles $g+c_j$ remain in addition;
original high external arrays keep their constant $c_j$ and acquire
$g(1/2)=0$. This gives a feasible state of the original complete
tree with at least one pair sum equal to one. That finite union of
closed boundary faces misses the unique strict maximum. Compactness
therefore supplies a fixed strict gap. This argument also treats a
single removed hub with only external arms and all closed endpoints.
